\section{What the Fathers and Doctors Actually Saw in the Cloths}
\label{sec:patristic}

Readers who turn from modern Easter anecdotes to the Fathers of the Church meet an immediate surprise. Far from speaking in a flat unison about John 20:3--10, the great Catholic commentators wrestle openly with the tension between verse 8 (``and he saw and believed'') and verse 9 (``For as yet they knew not the scripture, that he must rise again from the dead''). Two Doctors of the Church---St. John Chrysostom in the East and St. Augustine in the West---arrive at directly opposite conclusions about what the beloved disciple believed inside the tomb, while St. Gregory the Great, Blessed Theophylact of Ohrid, St. Thomas Aquinas, and Cornelius a Lapide navigate that very disagreement with remarkable candor and theological depth.

\subsection{St. John Chrysostom: Myrrh Like Lead and Calm Order}

Preaching in Antioch around the end of the fourth century, St. John Chrysostom saw in John 20:5--8 a direct refutation of the rumor spread by the chief priests and soldiers in Matthew 28:11--15---that the disciples came by night and stole the body. In \emph{Homily 85 on John} (\S 4), Chrysostom joins John 20:6--7 directly back to Nicodemus's hundred pounds of myrrh and aloes in John 19:39:
\begin{quote}
``When then she came and said these things, they hearing them, draw near with great eagerness to the sepulcher, and see the linen clothes lying, which was a sign of the Resurrection. For neither, if any persons had removed the body, would they before doing so have stripped it; nor if any had stolen it, would they have taken the trouble to remove the napkin, and roll it up, and lay it in a place by itself; but how? they would have taken the body as it was. On this account John tells us by anticipation that it was buried with much myrrh, which glues linen to the body not less firmly than lead; in order that when thou hearest that the napkins lay apart, thou mayest not endure those who say that He was stolen. For a thief would not have been so foolish as to spend so much trouble on a superfluous matter. For why should he undo the clothes? and how could he have escaped detection if he had done so? since he would probably have spent much time in so doing, and be found out by delaying and loitering. But why do the clothes lie apart, while the napkin was wrapped together by itself? That thou mayest learn that it was not the action of men in confusion or haste, the placing some in one place, some in another, and the wrapping them together.''
\end{quote}

Chrysostom does not stop at forensic logic. He notices the human portraits of the two Apostles and the humility of the Evangelist who records Peter's thoroughness ahead of his own:
\begin{quote}
``Observe too here again the absence of boastfulness in the Evangelist, how he witnesses to the exactness of Peter's search. For he himself having gotten before Peter, and having seen the linen clothes, inquired not farther, but withdrew; but that fervent one passing farther in, looked at everything carefully, and saw somewhat more, and then the other too was summoned to the sight. For he entering after Peter, saw the grave-clothes lying, and separate. Now to separate, and to place one thing by itself, and another, after rolling it up, by itself, was the act of some one doing things carefully, and not in a chance way, as if disturbed.''
\end{quote}

On the meaning of ``believed,'' Chrysostom is unequivocal across three consecutive homilies. In \emph{Homily 85.4} he states of both Apostles: ``From this they believed in the Resurrection. On this account Christ afterwards appeared to them, when they were convinced by what they had seen.'' In \emph{Homily 86.1}, explaining why Mary Magdalen lingered weeping while Peter did not, he writes that ``they having seen the linen clothes and believed, departed to their own homes in astonishment.'' And in \emph{Homily 87}, contrasting the two Apostles with Thomas's later demand to touch the wounds, he remarks that ``from the proof of the napkins, they straightway received the word concerning the Resurrection, and before they saw the body, exhibited all faith.''

\subsection{St. Cyril of Alexandria and St. Ambrose: Legal Witnesses and Easter Faith}

In fifth-century Alexandria, St. Cyril likewise read the linen cloths as the evidence that brought both Apostles to Resurrection faith, framing their race to the tomb through the Mosaic requirement of two witnesses (Deut 19:15). In Book XII of his \emph{Commentary on John}, Cyril writes:
\begin{quote}
``When these men (I mean Peter, and John the writer of this book, for he gives himself the name of the other disciple) heard these tidings from the woman's mouth, they ran with all the speed they could, and came to the sepulchre in haste, and saw the marvel with their own eyes, being in themselves competent to testify to the event, for they were two in number, as the Law enjoined. As yet they did not meet Christ risen from the dead, but infer His Resurrection from the bundle of linen clothes, and henceforth believed that He had burst asunder the bonds of death, as Holy Writ had long ago proclaimed that He would do. When, therefore, they looked at the issues of events in the light of the prophecies which turned out true, their faith was henceforth rooted on a firm basis.''
\end{quote}
One detail in Cyril's exposition deserves exact attention: Cyril speaks only of ``the bundle of linen clothes'' and never mentions the \latin{soudarion} separately in this passage. For him, the discarded burial bundle as a whole sufficed for the two Apostles to infer that the Lord had broken the bonds of death.

In the Latin West, St. Ambrose of Milan took a similarly positive view of both Apostles' faith. Commenting on Luke 24 in his \emph{Exposition of the Gospel according to Luke} (X.179), Ambrose declares that ``led on by so many examples of power, we believe that Peter could not doubt. It is plain that John also believed when he saw the Saviour---he who already then believed, after he saw the sepulchre empty of the body.''\endnote{St. Ambrose, \emph{Expositio evangelii secundum Lucam} X.179 (Migne, PL 15): \latin{``Tot igitur adducti virtutis exemplis Petrum credimus non potuisse dubitare. Joannem quoque credidisse manifestum est, cum Salvatorem videret: qui jam tunc credidit, posteaquam vacuum corporis vidit sepulcrum.''}}

\subsection{St. Augustine: The Force of Verse 9}

St. Augustine of Hippo read the exact same verses and stopped sharply at the word ``For'' (\latin{enim}, Greek \latin{gar}) opening verse 9. In \emph{Tractate 120 on the Gospel of John} (\S 9), preached to his congregation in Hippo, Augustine directly challenges the Resurrection-faith interpretation:
\begin{quote}
``And he saw, and believed. Here some, by not giving due attention, suppose that John believed that Jesus had risen again; but there is no indication of this from the words that follow. For what does he mean by immediately adding, \emph{For as yet they knew not the scripture, that He must rise again from the dead}? He could not then have believed that He had risen again, when he did not know that it behoved Him to rise again. What then did he see? what was it that he believed? What but this, that he saw the sepulchre empty, and believed what the woman had said, that He had been taken away from the tomb?''\endnote{St. Augustine, \emph{In Iohannis Evangelium Tractatus} 120.9 (Latin text: \latin{``Hic nonnulli parum attendentes, putant hoc Ioannem credidisse, quod Iesus resurrexit; sed quod sequitur, hoc non indicat. Quid sibi enim vult quod statim adiunxit: Nondum enim sciebant Scripturam, quia oportet eum a mortuis resurgere? Non ergo eum credidit resurrexisse, quem nesciebat oportere resurgere. Quid ergo vidit? quid credidit? Vidit scilicet inane monumentum, et credidit quod dixerat mulier, eum de monumento esse sublatum.''}). English from NPNF I.7 (trans. John Gibb and James Innes, 1888).}
\end{quote}

Augustine repeats this judgment consistently across his writings---in the opening of the very next sermon (\emph{Tractate 121.1}: ``they, when they came thither, found only the linen clothes wherewith the body had been shrouded; and what else could they believe but what she had told them, and what she had herself also believed?'') and in \emph{The Harmony of the Gospels} (III.24.69: ``Then John entered also, and saw in like manner, and believed what Mary had told him, namely, that the Lord had been taken away from the sepulchre''). To the obvious objection that Jesus had predicted His Resurrection plainly before His Passion, Augustine answers in \emph{Tractate 120.9} that the disciples were so accustomed to hearing the Lord speak in parables that they did not understand His plain speech and supposed He meant something else.

Augustine shows equal restraint in the preceding paragraph (\emph{Tractate 120.8--9}). Quoting Peter's sight of the linen cloths and the separate head-cloth, Augustine asks his hearers: ``Do we suppose these things have no meaning? I can suppose no such thing.'' Yet he refuses to invent an explanation on the spot, coining one of his brisk Latin wordplays on John arriving first and entering second: \latin{Prior venit, et posterior intravit. Neque hoc utique vacat, sed mihi ad ista non vacat}---``He came first, and entered last. Neither is this surely empty of meaning, yet I have not the leisure for these matters.''

\subsection{St. Gregory the Great: The Head-Cloth Rolled Up Without Beginning or End}

Nearly two centuries after Augustine, Pope St. Gregory the Great preached \emph{Homily 22 on the Gospels} in Rome on Easter Saturday (\latin{Sabbato in albis}), taking up the very spiritual mystery of the \latin{sudarium} that Augustine had left untouched. Gregory's contemplation of John 20:7 moves across three luminous dimensions (\emph{Homiliae in Evangelia} 22.3--4):

First, joining John 20:7 to St. Paul's declaration that ``the head of Christ is God'' (1 Cor 11:3), Gregory contemplates why the head-cloth lies apart from the body-linens:
\begin{quote}
``Why do we believe, brethren, that the head-cloth of the Lord is not found with the linen cloths in the sepulchre, unless because, as Paul attests, the head of Christ is God, and the incomprehensible mysteries of the Godhead are set apart from the knowledge of our weakness, and his power transcends the nature of creation?''
\end{quote}

Second, Gregory turns directly to the word \latin{involutum} (``rolled up'') and the phrase \latin{in unum locum} (``in one place''):
\begin{quote}
``And it is to be noted that it is said to be found not only apart, but also rolled up in one place. For of a linen cloth that is rolled up, neither its beginning nor its end is seen. Rightly therefore was the head-cloth found rolled up, because the height of the Godhead neither began to be nor ceases; neither is it born through a beginning nor straitened by an end. Well, moreover, is it added: \emph{In one place}, because God is not in a division of minds.''\endnote{St. Gregory the Great, \emph{Homiliae in Evangelia} 22.3--4 (Migne, PL 76): \latin{``Linteum quippe quod involvitur, ejus nec initium nec finis aspicitur. Recte ergo sudarium capitis involutum inventum est, quia celsitudo divinitatis nec coepit esse, nec desinit; nec per initium nascitur, nec termino coangustatur. Bene autem additur: In unum locum, quia in scissura mentium Deus non est.''} Here, as in Gregory's preceding and following sentences, the English rendering translates Migne's Latin plainly.}
\end{quote}

Third, drawing on the literal sense of \latin{sudarium} as a cloth that wipes the sweat of labor (a detail St. Thomas Aquinas highlights when excerpting Gregory in the \emph{Catena Aurea}), Gregory sees in the separated head-cloth the unique purity of Christ's Passion: ``Therefore the \latin{sudarium} that had been upon his head is found apart, because the Passion itself of our Redeemer is far set apart from our suffering, since he bore without guilt what we endure with guilt. He willed of his own accord to submit to death, to which we come unwillingly.'' In the same homily (excerpted in Aquinas's \emph{Catena Aurea} on John 20:1--9), Gregory also reads the race of the two Apostles ecclesially: John, the younger who reaches the tomb first yet does not enter until after Peter, figures the Synagogue, which knew the Law and prophecies yet hesitated before the mystery of the dead and risen Christ until the Gentile Church---figured by Peter, the elder---entered first, after whom Israel too shall enter in the fullness of time.

Yet what makes Gregory's homily so instructive for scriptural method is what comes in the very next paragraph (\emph{Homiliae in Evangelia} 22.5). Immediately after rising to the heights of mystical theology over the rolled-up \latin{sudarium}, Gregory turns to the literal meaning of ``and he saw, and believed''---and sides squarely with Augustine:
\begin{quote}
``What, brethren, what must he be judged to have believed? Was it that the Lord had risen, whom he was seeking? By no means [\latin{non utique}], because there was still darkness at the sepulchre, and the words that follow also contradict it when it is said: \emph{For as yet they knew not the Scriptures, that he must rise again from the dead}. What then did he see, and what did he believe? He saw the linen cloths lying, and believed what the woman had said, that the Lord had been taken away from the sepulchre.''
\end{quote}
Gregory felt no clash between sober literal restraint on verse 8 and rich contemplative theology on verse 7. He did not invent a first-century custom to justify his spiritual interpretation; he grounded his contemplation in the exact Latin words \latin{separatim}, \latin{involutum}, \latin{in unum locum}, and \latin{sudarium}.

\subsection{Blessed Theophylact of Ohrid: Chrysostom and Augustine Under One Roof}

A surprising bridge between East and West appears in the eleventh-century Byzantine archbishop Blessed Theophylact of Ohrid, whose \emph{Enarratio in Evangelium Ioannis} frequently condenses Chrysostom's homilies. Reading Theophylact directly at John 20 reveals a striking split.\endnote{Theophylact's commentary on John 20 is read here in Migne's facing Latin column (\emph{Patrologia Graeca} 124, cols. 289--291), as digitized from the 1864 printing; the parallel Greek column was not collated, and four minor character-recognition repairs in the Latin text are recorded in Appendix~\ref{app:scope}.}

On verses 6--7, Theophylact reproduces Chrysostom's entire anti-theft argument, even strengthening the comparison of myrrh: ``for this reason too the Evangelist said earlier that he was buried with much myrrh, which makes linen cling to the body no less than pitch [\latin{non minus bitumine}]; so that when we hear that the cloths were laid separately there, we may not believe those who say he was carried off by stealth. For no thief would have been so foolish as to spend so much care and time on a needless task.'' He likewise praises the Evangelist's freedom from vanity in attesting Peter's careful search, and adds a moral typology preserved in the \emph{Catena Aurea}: Peter represents the active life, prompt and fervent, while John represents the contemplative life, swift in intelligence yet outstripped in entering the mystery by Peter's ardent zeal.

Yet when Theophylact reaches verse 8---``and he saw, and believed''---he breaks sharply with Chrysostom and adopts Augustine's interpretation:
\begin{quote}
``He saw the burial cloths laid apart, and believed: not that the Lord had risen, but that he had been carried off by theft. For he believed Mary saying, \emph{They have taken away the Lord}. Why did he believe Mary and not think that he had risen? \emph{For as yet they had not known the Scripture, that he must rise again from the dead}, and therefore they believed Mary, and suspect theft and removal.''
\end{quote}
The divide over John 20:8 was never a simple geographic line between Greek East and Latin West; it lived inside the wording of verses 8 and 9 themselves, and even a Byzantine heir of Chrysostom felt the full grammatical pressure of verse 9.

\subsection{St. Thomas Aquinas and Cornelius a Lapide: The Scholastic Weighing}

In the thirteenth century, St. Thomas Aquinas brought the patristic conversation into clear scholastic focus in his \emph{Lectures on the Gospel of St. John} (cap. 20, lect. 1). First, asking why St. John narrates the placement of the cloths so carefully, Aquinas endorses Chrysostom's anti-theft proof---adding a realistic detail about the Roman guard from Matthew 27:66:
\begin{quote}
``To this it must be said, according to Chrysostom, that it was to exclude the false opinion spread abroad by the Jews, that the body of Christ had been stealthily removed. For if any had transferred it in that way, they would surely not have stripped him, especially since they would have had to do this hurriedly because of the presence of the guards; nor would they have been so careful to take up the \latin{sudarium}, and roll it up, and set it apart in one place, but would simply have left the cloths as they fell after taking the body. For this reason also he permitted himself to be buried with myrrh and aloes, which glue the linen cloths to the body, so that they could not so quickly be separated from the body.''\endnote{St. Thomas Aquinas, \emph{Super Evangelium S. Ioannis Lectura}, cap. 20, lect. 1 (Corpus Thomisticum text, Turin 1952 base). The clause about the guards (\latin{propter praesentiam custodum}) does not stand in the NPNF text of Chrysostom's homily; a parallel detail regarding the guards appears in Theophylact's Lucan exegesis.}
\end{quote}

Then, arriving at ``and he saw, and believed,'' Aquinas refuses to suppress either Doctor of the Church. He sets Augustine first as a formidable objection to the surface sense, and then shows how Chrysostom's exposition answers verse 9 syntactically:
\begin{quote}
``On first sight it could have this sense, that he saw those things and believed that Christ had risen. But, according to Augustine, this cannot stand, because there follows \emph{For as yet they knew not the Scripture}, and so forth. Therefore it must be said that he saw the empty tomb and believed what the woman had said, namely that someone had carried the Lord away. \dots\ Or, according to Chrysostom, he saw the linen cloths placed in such order as would not have been done if the body had been stealthily taken away, and he believed---with true faith---that Christ had risen from death. And then what follows, \emph{For as yet they knew not the Scripture}, refers back to what he says, \emph{he saw and believed}, as if to say: before he saw, he did not know the Scripture, that he must rise from the dead; but when he saw, he believed that he had risen from the dead.''
\end{quote}

Aquinas's grammatical solution for Chrysostom's side turns on the Greek pluperfect (\latin{\={e}deisan}, ``they had not yet understood''): verse 9 explains why the beloved disciple needed the physical evidence of the ordered cloths to come to faith, because up to that moment neither Apostle had yet understood from the Old Testament prophecies that the Messiah must rise from the dead, so only by stepping inside and seeing the tomb did John's faith awaken.

Four centuries later, the Jesuit exegete Cornelius a Lapide surveyed both camps in his \emph{Great Commentary} on John 20:8--10---listing St. Augustine, Theophylact, and Jansen on the taken-body side, and St. Cyril, St. Chrysostom, Euthymius, and St. Gregory of Nyssa on the Resurrection-faith side---and proposed a middle distinction that respects the singular verb of verse 8 alongside Luke 24:12:
\begin{quote}
``But this word `believed' more clearly and correctly applies only to S. John, who remembered the words of Christ, that He would rise on the third day. But Peter, on account of the strangeness of a Resurrection, and from his earnest desire to see Him alive again, was more slow to believe that Christ had risen. Whence the Angel significantly said to the women, `Go, tell His disciples and Peter' (Mark xvi. 7).''
\end{quote}
Summarizing verse 10 (``Then the disciples went away again unto their own home''), Lapide distills the two Apostles into four words: ``Peter wondering, John believing He had risen.''

\Needspace{10\baselineskip}
\begin{fourstable}{Commentator (Date)}{Significance of the Cloths}{Meaning of ``Believed'' (John 20:8)}{Notable Theological Accent}
St. John Chrysostom (late 4th c.) & Myrrh glues linen ``like lead''; ordered, rolled-up \latin{soudarion} refutes theft or hasty removal & Both Peter and John believed in the Resurrection from the cloths & John's humility in recording Peter's thoroughness; faith before seeing the risen body \\
\midrule
St. Ambrose (late 4th c.) & Empty sepulchre confirms Christ's prior displays of divine power & Both Peter (not doubting) and John (at the empty tomb) believed & Continuity between Christ's miracles in life and His Resurrection \\
\midrule
St. Augustine (early 5th c.) & Details surely carry meaning (\latin{neque hoc utique vacat}), left for another time; reconciles \latin{sind\={o}n}/\latin{othonia} & John believed Mary's report that the body had been taken away (governed by v. 9) & Sober grammatical rigor; reconciles Luke 24:12 and John 20:5--8 \\
\midrule
St. Cyril of Alexandria (early 5th c.) & ``Bundle of linen clothes'' (no separate mention of \latin{soudarion}) shows bonds of death broken & Both Apostles inferred the Resurrection from the linen cloths & Peter and John as two competent witnesses required by the Law (Deut 19:15) \\
\midrule
St. Gregory the Great (late 6th c.) & Separated, rolled-up \latin{sudarium} figures God's transcendent eternity, unity, and Christ's guiltless Passion & Follows Augustine (\latin{non utique}): John believed Mary's report that the Lord was taken away & Synagogue (John) and Gentile Church (Peter) running to the tomb \\
\midrule
St. Bede the Venerable (early 8th c.) & Reconciles Luke 24:12 and John 20:5--7: Peter stooped first, then entered before John & Focuses on Luke 24:12 (Peter departing in wonder) & Harmonizes Synoptic and Johannine narrative sequences \\
\midrule
Bl. Theophylact of Ohrid (late 11th c.) & Reproduces Chrysostom's anti-theft proof (myrrh adhering like pitch, \latin{bitumine}) & Breaks with Chrysostom to follow Augustine: believed Mary's report of theft & Active (Peter) and contemplative (John) lives approaching the mystery \\
\midrule
St. Thomas Aquinas (13th c.) & Endorses Chrysostom's anti-theft proof (adding haste from presence of guards) & Preserves both Augustine and Chrysostom side by side; explains v. 9 pluperfectly for Chrysostom & Balanced scholastic synthesis honoring both Doctors without flattening disagreement \\
\midrule
Cornelius a Lapide (early 17th c.) & Follows Chrysostom's anti-theft proof; summarizes earlier patristic consensus & Applies ``believed'' (Resurrection) specifically to John alone; Peter departs wondering & Joins John 20:8--10 to Luke 24:12 and the angel's ``and Peter'' (Mark 16:7) \\
\end{fourstable}

Two conclusions emerge cleanly from this long chorus of Catholic exegesis. First, the Church leaves the exegetical question of John 20:8 open: Catholics may follow Augustine and Gregory in understanding verse 8 as belief in Mary's report, or follow Chrysostom, Cyril, and Cornelius a Lapide in understanding it as the first awakening of Easter faith in the beloved disciple, with Aquinas holding both paths open. Second, across thirteen centuries of uninterrupted attention to every syllable of John 20:6--7, not one Father, Doctor, or commentator saw a dinner napkin---they saw the physical refusal of grave-robbery and the quiet sign of God's transcendent victory over death.
